\documentclass{beamer}
\setbeamertemplate{navigation symbols}{}

\usetheme{Copenhagen}
\useoutertheme{miniframes}

\makeatletter
\setbeamertemplate{headline}{%
  \begin{beamercolorbox}[wd=\paperwidth]{section in head/foot}
    \vskip2pt%
    \insertnavigation{\paperwidth}%
    \vskip2pt%
  \end{beamercolorbox}%
}
\makeatother

\usepackage{xcolor}
\usepackage{adjustbox} 
\usepackage{graphicx}
\usepackage{array}
\usepackage{xurl}
% \usepackage{tikz}
\usepackage{ulem}
\usepackage{minted}

\usepackage{polyglossia}

\setdefaultlanguage{czech}


\definecolor{FITBlue}{RGB}{0,115,188}
\definecolor{FITYellow}{RGB}{249,176,0}

\setbeamercolor{structure}{fg=FITBlue}
\setbeamercolor{palette secondary}{bg=FITBlue}
\setbeamercolor{palette tertiary}{fg=white,bg=FITYellow}


\setbeamertemplate{itemize item}{\raise1pt\hbox{\textbullet}}
\setbeamertemplate{itemize subitem}{\raise1pt\hbox{\textopenbullet}}

\setbeamercolor{itemize item}{fg=FITBlue}
\setbeamercolor{itemize subitem}{fg=gray!80!black}

\setbeamertemplate{enumerate item}{%
  \tikz[baseline=(n.base)]\node[fill=FITBlue, text=white, circle, inner sep=1pt, minimum size=1.2em](n){\footnotesize\bfseries\insertenumlabel};}

\setbeamertemplate{enumerate subitem}{%
  \tikz[baseline=(n.base)]\node[fill=gray!80!black, text=white, circle, inner sep=1pt, minimum size=1em](n){\footnotesize\insertenumlabel};}


\makeatletter
\newlength\beamerleftmargin
\setlength\beamerleftmargin{\Gm@lmargin}
\makeatother


\title{PA1 Konzultace}
\author{Klub FIT++}
\date{20. listopadu 2025}
\begin{document}

\maketitle

\section{Pointery}

\begin{frame}{Co jsou to pointery?}

\begin{itemize}
    \item \sout{jedničky a nuly, hvězdička} Datový typ
    \item \texttt{int*, double*, struct my\_data\_t*, long long**}
    \item Pointer aritmetika (+ 1 přičítá víc než jen 1)
    \item Speciální operace "dereference" \texttt{*variable, variable->xyz}
\end{itemize}
    
\end{frame}

\begin{frame}{Reprezentace pointerů}

\begin{itemize}
    \item "odkaz na proměnnou"
    \item adresa kde doopravdy proměnnou hledat
    \item adresa $0$ (\texttt{NULL}) se vnímá jako "nic zde není"
    \item pozor na na neplatné odkazy (neinicializované proměnné, zastaralá hodnota)
\end{itemize}

\end{frame}

\begin{frame}[fragile]{Jak s pointery pracovat?}

\begin{minted}{C}

int main() {
    int x = 1234;

    int *ptr = &x;

    printf("%d", *ptr);
}

\end{minted}
    
\end{frame}

\begin{frame}[fragile]{Praktické použití pointerů}

{\small
\begin{minted}[baselinestretch=0.8]{C}

int load_point(int *x, int*y);

int main() {
    int x = 0;
    int y = 0;

    if (load_point(&x, &y) != 0) {
        printf("Chybný vstup.\n");
        return;
    }
}

int load_point(int *x, int*y) {
    if (!x || !y) // if x or y is invalid pointer
        return 1;
    char c = 0;

    ...
}

\end{minted}
}

\end{frame}

\begin{frame}[fragile]

{\small
\begin{minted}[baselinestretch=0.8]{C}

int load_point(int *x, int*y) {
    if (!x || !y) // if x or y is invalid pointer
        return 1;
    char c = 0;

    int res = scanf(" [ %d , %d %c ", x, y, &c);

    if (*x < 0)
        *x = (-1) * (*x);
        
    if (*y < 0)
        *y = (-1) * (*y);

    // Error
    if (c != '[' || res != 3)
        return 2;

    return 0;
}
\end{minted}
}


\end{frame}

\begin{frame}[fragile]{Neplatné pointery}

\begin{minted}{C}

int load_point(int *x, int*y) {
    if (!x || !y) // if x or y is invalid pointer
        return 1;
    ...
}

int main() {
    int *x; // x is uninitialized => random address
    int *y; // x will be truthy, but *x throws error

    if (load_point(&x, &y) != 0) {
        printf("Chybný vstup.\n");
        return;
    }
}

\end{minted}

\end{frame}

\begin{frame}[fragile]{Neplatné pointery}

\begin{minted}{C}

int *get_constant() {
    int abc = 123;

    // We return pointer, but it outlives the variable
    return &abc;
}

int main() {
    int *abc = get_constant();

    // abc has correct value but might
    // already be overwritten by something else
    printf("%d\n", *abc);
}
\end{minted}

\end{frame}

\section{Struktury}

\begin{frame}{Struktury}

\begin{itemize}
    \item Často naše proměnná vyžaduje "kontext"
    \item Logické uspořádání více proměnných pod jednu
    \item "Uživatel", "Bod", "Vektor", jiná datová struktura - Halda, AVL strom...
\end{itemize}
    
\end{frame}

\begin{frame}[fragile]

{\small
\begin{minted}[baselinestretch=0.8]{C}

char nth_character(
  char str[],
  size_t length,
  size_t index) {
    if (index > length)
        return -1;

    return char[index];
}

int vec_push_int(
  int (*vec)[],
  size_t *size,
  size_t *capacity,
  int value) {
    if (*size >= *capacity) { ... }
    (*vec)[*size] = value;
    *size += 1;

    return 0;
}

\end{minted}
}

\end{frame}

\begin{frame}[fragile]{Příklad použití}

{\small
\begin{minted}[baselinestretch=0.8]{C}

typedef struct point_t {
    int m_x;
    int m_y;
} point_t;

void example_fn(point_t pt) {
    // work with pt.m_x, pt.m_y...
}

int main() {
    point_t mypoint;
    mypoint.m_x = 1;
    mypoint.m_y = 2;

    example_fn(mypoint);
}
    
\end{minted}
}
    
\end{frame}

\begin{frame}[fragile]{Pozor na inicializaci}

{\small
\begin{minted}[baselinestretch=0.8]{C}

typedef struct point_t {
    int m_x;
    int m_y;
} point_t;

void main() {
    point_t mypoint;

    if (mypoint.m_x > 0) { // Maybe true
        // ...
    }
}

typedef struct user_t {
    size_t m_id;
    // can be initialized to random address...
    char p_username[];
} user_t;
    
\end{minted}
}

\end{frame}

\begin{frame}[fragile]{Nejoptimálnější využívání}

{\small
\begin{minted}[baselinestretch=0.8]{C}

typedef struct vec_int_t {
  int *p_data;
  size_t m_size;
  size_t m_capacity;
}

int init_vec_int(vec_int_t *vec) {
    if (!vec)
        return -1;
    vec->m_size = 0;
    vec->m_capacity = 8;
    
    vec->p_data = (int *) malloc(...);
    if (!vec->p_data) {
      vec->m_capacity = 0;
      return -1;
    }

    return 0;
}

\end{minted}
}

\end{frame}

\begin{frame}[fragile]

{\small
\begin{minted}[baselinestretch=0.8]{C}

void delete_vec_int(vec_int_t *vec) {
    if (!vec)
        return;

    free(vec->p_data);
    vec->p_data = 0;
    vec->m_size = 0;
    vec->m_capacity = 0;
}

int main() {
    vec_int_t vec;
    if (!init_vec_int(&vec)) {
        return -1; // add error message?
    }

    // safely use vec
    delete_vec_int(&vec);
}

\end{minted}
}
\end{frame}

\section{Dynamická alokace}

\begin{frame}[fragile]{Zpátky k pointerům}

{\small
\begin{minted}[baselinestretch=0.8]{C}

int *get_constant() {
    int abc = 123;

    // How to actually make it so that the
    // value can be returned (and still be a pointer)
    // value can be returned (and still be a pointer)
    return &abc;
}

int main() {
    int *abc = get_constant();

    // abc might already be overwritten by something else
    printf("%d\n", *abc);     
}
\end{minted}
}
    
\end{frame}

\begin{frame}[fragile]{Problém vyřešen}

{\small
\begin{minted}[baselinestretch=0.8]{C}

int *get_constant() {
    int *abc_ptr = (int *) malloc(sizeof(int));
    if (abc_ptr)
        *abc_ptr = 123;

    return abc_ptr;
}

int main() {
    int *abc = get_constant();

    if (!abc) { ... }

    printf("%d\n", *abc);

    // Teď ale musíme manuálně říct, kdy abc už nepotřebujeme
    free(abc);
}
    
\end{minted}
}
    
\end{frame}

\begin{frame}{malloc/free}

\begin{itemize}
    \item malloc - "žádáme systém o n bytů místa pro libovolnou proměnnou"
    \item Tato proměnná je ale "za" pointerem. tj. přežije return funkce
    \item Abychom nevyplýtvali paměť musíme proměnnou zase vrátit (free)
    \item Ale nesmíme vrátit stejnou proměnnou dvakrát.
\end{itemize}

\end{frame}

\begin{frame}[fragile]{Užitečný příklad}

{\small
\begin{minted}[baselinestretch=0.8]{C}

typedef struct int_arr_t {
    size_t m_size;
    size_t m_last;
    int *p_data;
} int_arr_t;

int create_int_arr(int_arr_t *arr, size_t size) {
    int *ptr = (int *) malloc(size * sizeof(int));
    if (!ptr)
        return -1;

    arr->m_size = size;
    arr->m_last = 0;
    arr->p_data = ptr;

    return 0;
}

\end{minted}
}
    
\end{frame}

\begin{frame}[fragile]{}

{\small
\begin{minted}[baselinestretch=0.8]{C}

typedef struct int_arr_t {
    size_t m_size;
    size_t m_last;
    int *p_data;
} int_arr_t;

void delete_int_arr(int_arr_t *arr) {
    if (!arr)
        return;

    arr->m_size = 0;
    arr->m_last = 0;
    free(arr->p_data);
    arr->p_data = NULL;
}

\end{minted}

}

\end{frame}

\end{document}
